\chapter{Литературен обзор}

Развитието на адаптивното учене като механизъм за формиране на очакванията в макроикономическите модели представлява резултат от дългогодишна еволюция - от най-ранните модели на адаптивни очаквания през средата на XX век до съвременните динамични стохастични модели за общо равновесие (DSGE модели). Както подчертават \textcite{Eusepi_Preston2023}, адаптивното учене отразява по-реалистичен процес на вземане на решения в среда на несигурност и структурна промяна, при който икономическите агенти не разполагат с пълна информация за модела на икономиката и постепенно актуализират убежденията си.

Съвременната теория на адаптивното учене е формализирана основно в работите на \textcite{Evans2003}. 
Авторите изследват взаимодействието между очаквания, базирани на адаптивно учене и дизайна на паричната политика. Използваният модел е стандартен неокейнсиански модел, включващ IS-крива, крива на Филипс, и уравнение за парична политика. Адаптивното учене се въвежда чрез рекурсивен метод на най-малките квадрати (RLS). 

В анализа са взети предвид както оптимални правила, получени чрез минимизация на функция на загубите, така и приближени таргетиращи правила (approximate targeting), предложени от \textcite{Mccallum_Nelson2004}. Моделът допуска, че централната банка не познава напълно структурата на икономиката и трябва да оценява параметрите в движение, като използва тези оценки при определянето на лихвената политика. 

Резултатите показват, че при определени параметри паричната политика не успява да осигури сходимост към равновесие с рационални очаквания. Оптималните правила, извлечени чрез минимизация на функцията на загубите, не винаги гарантират E-устойчивост, докато някои опростени таргетиращи правила постигат стабилност в среда на адаптивни очаквания. 

В други изследвания основен фокус представлява наблюдаваната инерция  на инфлацията в икономиката. \textcite{MILANI20072065} показва, че чрез въвеждане на constant-gain учене е възможно да се възпроизведат нивата на инерция, наблюдавани в реалните икономически данни, без да се налага добавяне на механизми като навици или индексация на цените към предходни нива на инфлация в кривата на Филипс. В сходен дух \textcite{SLOBODYAN201226} интегрират адаптивно учене в DSGE модел на \textcite{Smets_Wouters2007} и достигат до извода, че моделът с учене демонстрира по-голяма инертност в симулираните времеви редове и предлага по-добро съвпадение с данните в сравнение с модела при рационални очаквания.

Класическият модел на \textcite{Gali_Monacelli2005} служи като основа за редица последващи разработки. На този фон \textcite{BaskSelander2007} изследват малка отворена икономика с хетерогенни очаквания, при която част от валутните търговци се ръководят от фундаментални фактори, а друга част - от chartist поведение. Chartism, в контекста тази статия, описва тип техническа търговия, при която търговците базират своите очаквания за бъдещия обменен курс върху минали тенденции и движения в цената, а не върху фундаментални икономически фактори. 

Основната цел на изследването е да се оцени как включването на обменния курс в правилото на Тейлър, в комбинация с хетерогенни очаквания на валутния пазар, влияе върху равновесието при адаптивно учене чрез рекурсивен метод на най-малките квадрати (RLS).

Авторите показват, че наличието на умерен дял chartist-търговци може да промени границите на E-устойчивост, както и че реакцията на паричната политика спрямо обменния курс може да действа стабилизиращо, ако е умерена. Прекалено силна реакция води до загуба на устойчивост.

\textcite{Chen2008} допълват анализа за реакцията на централната банка. Авторите анализират четири алтернативни правила на паричната политика в контекста на японската икономика и адаптивното формиране на очаквания. Базовото правило представлява стандартна версия на правилото на Тейлър, с коефициенти 1.5 пред инфлацията и 0.5 пред отклонението на производството от потенциалната му стойност. Второто правило отразява по-„стриктна“ реакция към инфлацията  Третото включва стабилизация на условията на търговия като цел на паричната политика. Последното правило заменя CPI инфлацията с инфлация на вътрешните производствени цени (DPP). За всяко правило авторите изчисляват динамиката на инфлацията и отклонението на производството от потенциалната му стойност под три различни механизма на очаквания - рационални очаквания, рекурсивни най-малки квадрати (RLS) и constant-gain учене. За всяко правило резултатите се оценяват чрез функция на загубите, която обобщава влиянието на волатилността на инфлацията и отклонението на производството.

Резултатите показват, че адаптивното учене води до по-висока волатилност и съответно до по-големи загуби спрямо рационалните очаквания. Второ, по-агресивните правила за контрол на инфлацията могат да увеличат колебанията в производството, но въпреки това намаляват инфлационната нестабилност и намаляват общата функция на загубите.

Ролята на адаптивното учене е изследвана и в контекста на малки отворени икономики с по-изразена инфлационна динамика. \textcite{BASDEVANT20051074}  въвежда механизъм на адаптивно учене, чрез който икономическите участници постепенно актуализират своите вярвания за динамиката на инфлацията и реалната икономика въз основа на наличните данни и нова информация. Авторът допуска, че процесът на учене може да отслабва с времето, така че прогнозите постепенно да се приближат до рационалните и на практика да спрат да се променят. 

Целта на анализа е да се установи доколко наблюдаваната инерция на инфлацията в Нова Зеландия може да се обясни не толкова чрез структурни фактори, колкото чрез начина, по който икономическите агенти формират своите очаквания. За тази цел е използван малък структурен модел, включващ основните уравнения на неокейнсианската рамка.

Резултатите от симулациите показват, че динамиката на инфлацията в Нова Зеландия може да бъде възпроизведена дори без въвеждане на реални или номинални фрикции - единствено чрез адаптивен процес на очаквания. В допълнение при симулации с по-бърз темп на учене инфлацията реагира по-силно на промени в паричната политика и се стабилизира по-бързо, докато по-бавният процес на учене води до по-продължителна инерция.

В друга разработка от \textcite{DIZIOLI2024100119} се анализира защо централните банки в развиващите се икономики реагират по-рано и по-агресивно при очаквания за повишение на инфлацията. Авторите изследват този въпрос чрез малък DSGE модел, в който очакванията се формират чрез адаптивно учене. Моделът е оценен с данни за Бразилия и САЩ за периода 2000–2019 г.

Първият основен резултат е, че в модела адаптивното учене улавя по-реалистично динамиката на инфлацията. Адаптивните очаквания създават по-голяма инерция, като инфлационните очаквания реагират по-силно на минала инфлация и остават повишени за по-дълго време. Този механизъм е особено отчетлив в Бразилия, където се оказва, че очакванията са по-чувствителни към шокове, отколкото в САЩ.

Вторият ключов резултат е свързан с поведението на оптималната парична политика. Поради по-голямата инерция адаптивните очаквания се отдалечават по-силно от равновесната траектория и за по-дълго време от равновесната траектория. Това налага централната банка да повиши лихвения процент в по-голям мащаб.

 \textcite{Fukac2008} разглежда сценарий, в който централната банка и частният сектор формират различни очаквания за икономиката и учат с различни скорости. Получените резултати показват, че прекалено агресивната реакция на паричната политика към инфлацията увеличава краткосрочната волатилност, но в дългосрочен план адаптивното учене води до сближаване на очакванията и стабилизиране на икономиката. По този начин авторът подчертава, че степента на хетерогенност в очакванията има важно значение за динамиката на икономиката и оптималния дизайн на политиката.

\textcite{MUTO201152} поставя акцент върху взаимодействието между прогнозите на централната банка и очакванията на частните агенти. Докато стандартните модели приемат, че и двете групи независимо формират очакванията си чрез адаптивно учене, емпирични данни за Япония показват, че прогнозите на централната банка служат като ориентир за частните агенти. На тази основа Muto въвежда хетерогенност в механизма на възприеманата динамика: централната банка и частният сектор използват различни възприятия за икономическите зависимости, но прилагат една и съща алгоритмична форма на учене (RLS).

Авторът показва, че когато частните агенти изграждат своите очаквания, изхождайки от прогнозите на централната банка, грешките в техните очаквания стават по-големи от тези на самата централна банка. Причините са две: първо, частните агенти допускат грешки при оценяване на параметрите в собствените си възприемани правила на движение; и второ, всяка грешка на централната банка директно се пренася в техните очаквания. В резултат общата грешка в прогнозите на частния сектор значително превишава тази на централната банка. Поради това, когато централната банка включва собствените си прогнози в правилото на Тейлър, нейният стандартен отговор към очакваната инфлация се оказва недостатъчен за стабилизиране на икономиката. За да бъде постигната устойчивост на Е-равновесието, е необходимо централната банка да реагира значително по-силно на очакваната инфлация, отколкото предписва принципът на Тейлър.

Muto разглежда и обратната ситуация, при която централната банка е тази, която се учи, използвайки прогнозите на частния сектор. Тогава грешките в прогнозите на централната банка стават по-големи от тези на частните агенти, което води до огледален резултат: за осигуряване на устойчивост на Е-равновесието не е необходима реакция по-силна от тази, предписана от стандартния принцип на Тейлър. В този режим, когато очакванията на централната банка изхождат от тези на частния сектор, изискванията към паричната политика са по-слаби.

Прегледът на литературата показва, че приложението на адаптивното учене в неокейнсианските модели води до разнообразни резултати, като различните автори използват различни методологични подходи, структури на модела и допускания за формирането на очаквания. В едни изследвания паричната политика трябва да бъде значително по-агресивна, за да стабилизира процеса на учене, докато в други прекомерната реакция води до засилена краткосрочна волатилност. В някои модели адаптивното учене предлага по-добро съвпадение с емпиричните данни и обяснява наблюдаваната инерция, докато в други основният акцент пада върху взаимодействието между частния сектор и централната банка и ролята на информационните разлики.

Общото между всички разгледани изследвания е, че адаптивното учене променя фундаментално динамиката на модела, като влияе върху стабилността на равновесието, характерa на инфлационната инерция и оптималния дизайн на паричната политика. Независимо от разликите в методологията, литературата единодушно подчертава, че очакванията, формирани чрез процеси на учене, не са пасивен елемент, а ключов механизъм, който оформя както краткосрочните, така и дългосрочните макроикономически резултати.

 
